\documentclass[12pt]{article}
\usepackage[utf8]{inputenc}
\usepackage[T1]{fontenc}
\usepackage{hyperref}
\title{Dari Language Resources for K-12 Curriculum Development}
\author{Compiled by Claude Sonnet 5.5 (claude-sonnet-5-5) for Andr\'as Kornai}
\date{}
\begin{document}
\maketitle

\section{Introduction}
Dari (ISO 639-3 prs) is an official language of Afghanistan and a language of instruction. The search surfaced few primary sources: no official Ministry of Education curriculum documents appeared, and the third-party pages found describe grade-level syllabi in general terms. They should be checked against Ministry publications before use.

\section{Primary Educational Resources}
\begin{itemize}
\item Third-party summaries of the Dari curriculum for Grade 1 and Grade 9 (unofficial; detailed syllabi are said to be hard to obtain): \url{https://learner.net/u/afghanistan-grade-1-dari-language} and \url{https://learner.net/u/afghanistan-grade-9-dari-curriculum}
\item Overview of the lower secondary syllabus (Grades 7--9), also unofficial: \url{https://learner.net/u/afghanistan-grade-7-9-dari-studies}
\end{itemize}
The search summary also mentioned an ``Afghan Children Read'' collection of about 45 Dari early-reading books on Bloom Library, but no URL for it appeared in the results, so it is not listed with a link.

\section{Linguistic and Reference Materials}
No grammar or dictionary of Dari was found in these results. The Georgetown University Press adult-learner textbooks by Rahman Arman were mentioned, but only bookshop pages appeared, so they are not listed.

\section{Religious and Community Resources}
None found.

\section{Digital and Audio Resources}
None found with a verifiable URL.

\section{Curriculum Development Context}
Dari already has a national school curriculum. The main task for a curriculum project is obtaining the official textbooks and teacher guides from the Ministry of Education or reliable repositories, which this search could not locate.
\end{document}